The Frevia project adopts Scrum as its primary management approach, with a fixed delivery window of sixteen weeks divided into five sprints. Scrum is used because the product is not fully specified at the start: the payment, dispute, and verification flows each depend on decisions that only become clear once an earlier increment exists and has been reviewed. An iterative process allows the team to deliver a working slice of the marketplace in every sprint instead of accumulating unverified work until the end.

Three artefacts govern the work. The Product Backlog holds the 214 SMAR items listed in Section 1.1, each with an identifier, a complexity rating, and an effort estimate in man-days. The Sprint Backlog is the subset committed at Sprint Planning, and it is capped at 67 man-days for Sprints 2 to 5 and at 104 man-days for the four-week first sprint. The Increment is what the team demonstrates at the Sprint Review: a running system in which the committed items are usable end to end, not a set of unfinished branches.

The cadence differs deliberately between the first sprint and the rest. Sprint 1 runs for four weeks because it carries the architectural decisions that everything else depends on: the monorepo layout, the shared validation package, the data model, the authentication flow, and the permission guard. Sprints 2 to 5 run for three weeks each and share an identical effort envelope, which makes velocity comparable across them and makes an overrun visible early rather than at the end of the project. Each sprint ends with a review with the supervisor and a retrospective, and the retrospective produces at most two concrete process changes so that the list of improvements stays actionable.

Work is organised by functional module rather than by technical layer. An item such as dual contract signing therefore includes its schema change, its API endpoint, its permission entry, its interface, and its tests, and it is only counted as done when all of those parts exist. This keeps the increment demonstrable and prevents the common failure mode in which the backend is "finished" for three sprints while the interface that exercises it does not exist. The Definition of Done used for every item is stated in Section 2.1.

Communication is continuous rather than ceremonial. A short daily stand-up covers progress, the next item, and blockers; blockers that affect another member are resolved immediately after the stand-up rather than carried to the next day. Source code, the task board, and the sprint burndown live in the same GitHub repository that holds the code, so a review comment and the change it refers to are never more than one click apart. Decision records, meeting minutes, and design notes are kept in the shared documentation folder referenced in Section 6.

Risk is reviewed at every Sprint Planning rather than only at the start of the project. The register in Section 1.3 is re-scored when a new integration is introduced, and any risk whose possibility increases is paired with a concrete action in the sprint backlog instead of being left as a note. This is what turns the register into a management instrument rather than a document produced once for the report.

By applying Scrum with a fixed effort envelope, a module-oriented breakdown, and an explicit Definition of Done, the team keeps the product in a releasable state at the end of every sprint and can trade scope, rather than quality or schedule, when reality diverges from the plan.
